\documentclass[UTF8,a4paper,12pt]{ctexart}
\usepackage[margin=2.5cm]{geometry}
\usepackage{amsmath,amssymb,bm}
\usepackage{booktabs,multirow,array}
\usepackage{graphicx}
\usepackage{float}
\usepackage{siunitx}
\usepackage{hyperref}
\usepackage{enumitem}
\usepackage{xcolor}

\hypersetup{colorlinks=true,linkcolor=black,citecolor=black,urlcolor=blue}
\graphicspath{{figures/}}
\setlength{\parindent}{2em}
\setlength{\parskip}{0pt}

\title{ {{COMPETITION_TITLE}} }
\author{参赛队伍：\underline{\hspace{4cm}}}
\date{}

\begin{document}
\maketitle

\begin{abstract}
% 最后撰写。逐问说明任务、核心模型、验证方式和可追溯的主要数值。
[待填写]

\textbf{关键词：} [待填写]
\end{abstract}

\section{问题重述}
[待填写]

\section{问题分析}
% 逐问说明：任务要求、关键难点、数据/变量特征、模型选择理由、求解与检验方法。
% 明确前一问输出如何进入后一问，不重复原题，不在此处堆公式。
[待填写：跨问叙事与总体模型链]

\section{模型假设}
\begin{enumerate}[label=\arabic*.]
  \item [待填写：团队引入的假设或口径、理由、适用范围及潜在偏差]
\end{enumerate}

\section{符号说明}
\begin{table}[H]
  \centering
  \caption{主要符号}
  \begin{tabular}{lll}
    \toprule
    符号 & 含义 & 单位或范围 \\
    \midrule
    $x$ & [待填写] & [待填写] \\
    \bottomrule
  \end{tabular}
\end{table}

\section{数据预处理与审计}
[待填写：原始问题、处理方法、关键参数及依据、保留或删除的信息、处理效果证据]

\section{模型建立与求解}
\subsection{问题一}
\subsubsection{任务、难点与模型选择}
[待填写：为什么该模型适合本问；与上一问的承接或与下一问的依赖]

\subsubsection{基线模型与诊断}
[待填写：基线定义、结果、暴露出的失败模式]

\subsubsection{主模型的建立}
% 每个重要关系按“文字动机 → 公式 → 逐项解释 → 计算用途”展开。
[待填写]

\begin{equation}
  f(\mathbf{x}) = [待填写]
  \label{eq:q1-model}
\end{equation}
式中各符号含义、单位及式 \eqref{eq:q1-model} 的用途为：[待填写]。

\subsubsection{求解算法}
% 说明初始化、迭代更新、约束处理、停止条件；不要写成库函数调用日志。
[待填写]

\subsubsection{结果与验证}
[待填写：先写结论，再给图表/数值证据，用模型机制解释，最后说明限制或下游用途]

\subsubsection{本问小结与下一问衔接}
[待填写：直接答案、最强证据、传递给下一问的变量/参数/约束]

\subsection{后续问题}
% 按“任务与选择—基线—模型建立—求解—结果验证—小结衔接”复制问题一结构，
% 并明确继承了什么、增加了什么、为何需要改变模型。
[待填写]

\section{敏感性、稳健性与误差分析}
[待填写]

\section{模型评价与推广}
\subsection{优点}
[待填写：只写已经验证的性质，例如误差、可行性或稳定性]

\subsection{局限}
[待填写：说明影响范围、可能偏差方向和针对性改进]

\subsection{推广条件}
[待填写：新场景中保持不变的机制，以及必须重新标定的数据或参数]

\section{结论}
[待填写：按题目顺序回答，不引入新结果]

\begin{thebibliography}{99}
\bibitem{ref1} [待核实的真实文献]
\end{thebibliography}

\appendix
\section{详细推导}
[待填写：正文须引用本附录，公式与正文保持一致]

\section{核心代码与补充结果}
[待填写：代码应可运行并对应正文结果，删除无关调试内容]

\end{document}
